\section{La Meta (Goldratt) - Heurísticas V/F}

\begin{tcolorbox}[colback=blue!5!white,colframe=blue!75!black,title=Mapa Mental Rápido (Cópialo)]
\textbf{Objetivo:} Ganar dinero (Aumentar Throughput $\uparrow$, Bajar Inventario $\downarrow$, Bajar Gasto Operativo $\downarrow$).\\
\textbf{Cuello de Botella (CB):} Recurso con Capacidad $\le$ Demanda (Herbie).\\
\textbf{Regla de Oro:} Una hora ganada en un CB es una hora ganada para el sistema. Una hora ganada en un NO-CB es un espejismo.
\end{tcolorbox}

\subsection{Heurísticas Verdadero / Falso (Trampas Clásicas)}
Anota estas reglas. Si lees lo contrario en el examen, es \textbf{FALSO}.

\begin{itemize}[itemsep=1ex]
    \item \textbf{Sobre Tamaños de Lote:}
    \begin{itemize}
        \item Reducir el lote de transferencia $\rightarrow$ Disminuye el WIP (Inventario en Proceso) y el Lead Time (Tiempo de Flujo).
        \item Reducir el lote $\rightarrow$ \textbf{NO} aumenta el Lead Time (es falso si dice lo contrario).
        \item Los lotes de transferencia \textbf{no} tienen que ser iguales a los lotes de proceso.
    \end{itemize}
    
    \item \textbf{Sobre Cuellos de Botella (CB):}
    \begin{itemize}
        \item Optimizar un recurso no-CB es una pérdida de tiempo y solo genera más inventario (WIP) atascado frente al CB.
        \item El CB dicta el ritmo (Tambor / Drum) de toda la planta.
        \item Si se daña una máquina frente al CB, importa menos que si se daña el propio CB.
    \end{itemize}
    
    \item \textbf{Sobre las Fluctuaciones Estadísticas:}
    \begin{itemize}
        \item Las fluctuaciones y eventos dependientes se acumulan en lugar de promediarse (fenómeno de la excursión de boy scouts).
        \item El último paso de la línea siempre sufrirá los peores retrasos heredados de los procesos anteriores.
    \end{itemize}
\end{itemize}

\subsection{Los 5 Pasos de la Teoría de Restricciones (TOC)}
(Anota el orden exacto, es evaluado siempre):
\begin{enumerate}
    \item \textbf{Identificar} la restricción del sistema (Encontrar a Herbie).
    \item \textbf{Explotar} la restricción (Que Herbie no se detenga por tonterías, quitarle la carga inútil).
    \item \textbf{Subordinar} todo lo demás a la decisión anterior (Los rápidos no deben adelantar a Herbie; sistema de cuerdas).
    \item \textbf{Elevar} la restricción (Comprar otra máquina, contratar más gente para el CB).
    \item \textbf{Volver al paso 1} (Si se rompió la restricción, buscar el nuevo CB. No dejar que la inercia gane).
\end{enumerate}
